\subsection{Integration of positive functions}

Recall that we have agreed to define the product \(0 \cdot (+\infty)\) to be equal to 0. This convention has in particular the following consequence: if f is a numerical function \(\geqslant 0\) defined on a locally compact space equipped with a positive measure \(\lambda\) , then \(\lambda^{*}(af) = a \cdot \lambda^{*}(f)\) for every constant a such that \(0 \leqslant a \leqslant +\infty\) . This is obvious if a = 0; if \(a = +\infty\) , then \(\lambda^{*}(af) = a \cdot \lambda^{*}(f) = 0\) or \(\lambda^{*}(af) = a \cdot \lambda^{*}(f) = +\infty\) according as f is \(\lambda\) -negligible or not; finally, if \(0 < a < +\infty\) , one knows that \(\lambda^{*}(af) = a \cdot \lambda^{*}(f)\) .

PROPOSITION 5. --- Let f be a numerical function \(\geqslant 0\) , lower semicontinuous on \(T \times T'\) . Then, the function

\[
t \mapsto \int^ {*} f (t, t ^ {\prime}) d \mu^ {\prime} (t ^ {\prime})
\]

is lower semi-continuous on T, and

\[
\iint^ {*} f (t, t ^ {\prime}) d \mu (t) d \mu^ {\prime} (t ^ {\prime}) = \int^ {*} d \mu (t) \int^ {*} f (t, t ^ {\prime}) d \mu^ {\prime} (t ^ {\prime}).\tag{3}
\]

This is a consequence of Prop. 2 of §3, No.~1, taking into account Lemma 1 of No.~1.

COROLLARY 1. --- Let \(f\) (resp. \(f'\) ) be a lower semi-continuous function \(\geqslant 0\) defined on \(T\) (resp. \(T'\) ); the function \(f \otimes f': (t, t') \mapsto f(t)f'(t')\) is then lower semi-continuous on \(T \times T'\) , and

\[
\iint^ {*} f (t) f ^ {\prime} \left(t ^ {\prime}\right) d \mu (t) d \mu^ {\prime} \left(t ^ {\prime}\right) = \left(\int^ {*} f (t) d \mu (t)\right) \left(\int^ {*} f ^ {\prime} \left(t ^ {\prime}\right) d \mu^ {\prime} \left(t ^ {\prime}\right)\right).\tag{6}
\]

Let \(G\) (resp. \(G'\) ) be the set of functions \(g \in \mathcal{H}_{+}(T)\) (resp. \(g' \in \mathcal{H}_{+}(T')\) ) such that \(g \leqslant f\) (resp. \(g' \leqslant f'\) ); then

\[
f \otimes f ^ {\prime} = \sup _ {g \in \mathrm{G}, g ^ {\prime} \in \mathrm{G} ^ {\prime}} g \otimes g ^ {\prime}.
\]

Since the functions \(g \otimes g'\) belong to \(\mathcal{H}_{+}(\mathrm{T} \times \mathrm{T}')\) , \(f \otimes f'\) is indeed lower semi-continuous, and (6) follows at once from Prop. 5 (or directly by passage to the limit in the preceding formula).

COROLLARY 2. --- Let A be a \(\mu\) -moderated subset of T, and A' a \(\mu'\) -moderated subset of T'; then A × A' is \(\nu\) -moderated in T × T'.

In view of the definition of moderated set (§1, No.~2, Prop. 5), it suffices to show that if B is an integrable open set in T, and B' an integrable open set in T', then the open set B × B' is integrable. This follows at once from Cor. 1.

COROLLARY 3. --- Let A be a \(\mu\) -negligible subset of T, and let B' be a \(\mu'\) -moderated subset of T'; then A × B' is \(\nu\) -negligible.

For, \(A \times B'\) is locally \(\nu\) -negligible (Prop. 4) and \(\nu\) -moderated (Cor. 2), hence is \(\nu\) -negligible ( \(\S1\) , No.~2, Cor. 1 of Prop. 7).

Cors. 2 and 3 may be extended to the product of two complex measures, on applying the statement to their absolute values (Ch. III, §4, No.~2, Prop. 3).

PROPOSITION 6. --- Let \(f\) be a numerical function \(\geqslant 0\) defined on \(\mathbf{T} \times \mathbf{T}'\) . Then

\[
\iint^ {*} f (t, t ^ {\prime}) d \mu (t) d \mu^ {\prime} (t ^ {\prime}) \geqslant \int^ {*} d \mu (t) \int^ {*} f (t, t ^ {\prime}) d \mu^ {\prime} (t ^ {\prime}).
\]

This follows from Prop. 3 of §3, No.~2, on taking (2) into account.

PROPOSITION 7. --- Let \(f\) be a \(\nu\) -measurable positive numerical function defined on \(\mathbf{T} \times \mathbf{T}'\) .

\begin{enumerate}
\def\labelenumi{\alph{enumi})}
\tightlist
\item
  If \(f\) is \(\nu\) -moderated, then the functions \(t \mapsto \int^{*} f(t, t') d\mu'(t')\) , \(t' \mapsto \int^{*} f(t, t') d\mu(t)\) are measurable and moderated for \(\mu\) and \(\mu'\) , respectively, and
\end{enumerate}

\[
\begin{array}{r} \iint^ {*} f (t, t ^ {\prime}) d \mu (t) d \mu^ {\prime} (t ^ {\prime}) = \int^ {*} d \mu (t) \int^ {*} f (t, t ^ {\prime}) d \mu^ {\prime} (t ^ {\prime}) \\ = \int^ {*} d \mu^ {\prime} (t ^ {\prime}) \int^ {*} f (t, t ^ {\prime}) d \mu (t). \end{array}\tag{7}
\]

\begin{enumerate}
\def\labelenumi{\alph{enumi})}
\setcounter{enumi}{1}
\tightlist
\item
  If the measure \(\mu'\) is moderated, then the function \(t \mapsto \int^{\bullet} f(t, t') d\mu'(t')\) is \(\mu\) -measurable, and
\end{enumerate}

\[
\iint^ {\bullet} f (t, t ^ {\prime}) d \mu (t) d \mu^ {\prime} (t ^ {\prime}) = \int^ {\bullet} d \mu (t) \int^ {\bullet} f (t, t ^ {\prime}) d \mu^ {\prime} (t ^ {\prime}).\tag{8}
\]

The assertion a), as well as the assertion b) when \(\mu'\) is bounded, are consequences of Prop. 5 of §3, No.~2. To treat the case that \(\mu'\) is moderated, let us represent \(\mu'\) as a sum \(\sum_{n\in\mathbf{N}}\mu_n'\) of a sequence of bounded measures (§2, No.~3, Prop. 4). The function \(t\mapsto\int^{\bullet}f(t,t')d\mu_n'(t')\) is then \(\mu\) -measurable, and

\[
\iint^ {\bullet} f (t, t ^ {\prime}) d \mu (t) d \mu_ {n} ^ {\prime} (t ^ {\prime}) = \int^ {\bullet} d \mu (t) \int^ {\bullet} f (t, t ^ {\prime}) d \mu_ {n} ^ {\prime} (t ^ {\prime}).
\]

But \(\mu \otimes \mu' = \sum_{n \in \mathbf{N}} (\mu \otimes \mu_n')\) (Prop. 1); the assertion b) is then obtained by summing on \(n\) ( \(\S 2\) , No.~2, Prop. 1).

COROLLARY 1. --- Let H be a subset of T × T', and let A be the set of t ∈ T such that the section H(t) of H at t is not μ'-negligible.

\begin{enumerate}
\def\labelenumi{\alph{enumi})}
\item
  If H is \(\nu\) -negligible then A is \(\mu\) -negligible.
\item
  If H is locally \(\nu\) -negligible and \(\mu'\) is moderated, then A is locally \(\mu\) -negligible.
\end{enumerate}

The property a) follows at once from Prop. 7 (or from Prop. 6). Under the hypotheses of b), it comes to the same to say that \(H(t)\) is locally \(\mu'\) -negligible or that it is \(\mu'\) -negligible, since \(\mu'\) is moderated ( \(\S1\) , No.~2, Prop. 7). Thus property b) follows from the formula (8).

This corollary extends at once, by passage to absolute values, to the product of two complex measures. The same is true for the following corollary:

COROLLARY 2. --- If a set \(A \subset T \times T'\) is \(\nu\) -integrable, then the section \(\mathrm{A}(t)\) of A at t is \(\mu'\) -integrable for almost every \(t \in T\) , the function \(t \mapsto \mu'(\mathrm{A}(t))\) is \(\mu\) -integrable, and

\[
\nu (\mathrm{A}) = \int \mu^ {\prime} (\mathrm{A} (t)) d \mu (t).\tag{9}
\]

PROPOSITION 8. --- For every pair of numerical functions \(f \geqslant 0\) , \(f' \geqslant 0\) defined on T and T', respectively,

\[
\iint^ {\bullet} f (t) f ^ {\prime} \left(t ^ {\prime}\right) d \mu (t) d \mu^ {\prime} \left(t ^ {\prime}\right) = \left(\int^ {\bullet} f (t) d \mu (t)\right) \left(\int^ {\bullet} f ^ {\prime} \left(t ^ {\prime}\right) d \mu^ {\prime} \left(t ^ {\prime}\right)\right).\tag{10}
\]

We begin by treating the case that \(\mu\) and \(\mu'\) are measures with compact support; the same is then true of \(\mu \otimes \mu'\) , and all of the symbols \(\int^{\bullet}\) , \(\iint^{\bullet}\) may be replaced by upper integrals. By Prop. 6,

\[
\begin{array}{r l} \iint^ {*} f (t) f ^ {\prime} (t ^ {\prime}) d \mu (t) d \mu^ {\prime} (t ^ {\prime}) & \geqslant \int^ {*} d \mu (t) \int^ {*} f (t) f ^ {\prime} (t ^ {\prime}) d \mu^ {\prime} (t ^ {\prime}) \\ & = \left(\int^ {*} f (t) d \mu (t)\right) \left(\int^ {*} f ^ {\prime} (t ^ {\prime}) d \mu^ {\prime} (t ^ {\prime})\right). \end{array}
\]

To establish the reverse inequality, let us choose a function \(h \geqslant f\) (resp. \(h' \geqslant f'\) ), the lower envelope of a sequence \((h_n)\) (resp. \((h'_n)\) ) of lower semi-continuous functions, such that

\[
\int^ {*} h (t) d \mu (t) = \int^ {*} f (t) d \mu (t)
\]

\((\text{resp. } \int^{*} h'(t') d\mu'(t') = \int^{*} f'(t') d\mu'(t'))\) ; the existence of such functions follows immediately from the definition of upper integral (Ch. IV, §1, No.~3, Def. 3) and Lebesgue's theorem. Applying Prop. 7 to the measurable function \(h \otimes h'\) , we have

\[
\begin{array}{l} \iint^ {*} f (t) f ^ {\prime} (t ^ {\prime}) d \mu (t) d \mu^ {\prime} (t ^ {\prime}) \leqslant \iint^ {*} h (t) h ^ {\prime} (t ^ {\prime}) d \mu (t) d \mu^ {\prime} (t ^ {\prime}) \\ = \left(\int^ {*} h (t) d \mu (t)\right) \left(\int^ {*} h ^ {\prime} (t ^ {\prime}) d \mu^ {\prime} (t ^ {\prime})\right) \\ = \left(\int^ {*} f (t) d \mu (t)\right) \left(\int^ {*} f ^ {\prime} (t ^ {\prime}) d \mu^ {\prime} (t ^ {\prime})\right), \end{array}
\]

which is the sought-for inequality. Thus the proposition is established when \(\mu\) and \(\mu'\) are measures with compact support. To treat the general case, it suffices to represent \(\mu\) (resp. \(\mu'\) ) as the sum of a family \((\mu_{\alpha})_{\alpha\in\mathbb{A}}\) (resp. \((\mu'_{\beta})_{\beta\in\mathbb{B}}\) ) of measures with compact support ( \(\S2\) , No.~3, Prop. 4), write the formula (10) for each measure \(\mu_{\alpha}\otimes\mu'_{\beta}\) , and sum on \((\alpha,\beta)\) , taking into account Prop. 1 ( \(\S2\) , No.~2, Prop. 1).

COROLLARY 1. --- With notations as in Prop. 8,

\[
\iint^ {*} f (t) f ^ {\prime} \left(t ^ {\prime}\right) d \mu (t) d \mu^ {\prime} \left(t ^ {\prime}\right) = \left(\int^ {*} f (t) d \mu (t)\right) \left(\int^ {*} f ^ {\prime} \left(t ^ {\prime}\right) d \mu^ {\prime} \left(t ^ {\prime}\right)\right)\tag{11}
\]

except possibly when one of the factors of the second member is equal to 0 and the other is equal to +∞.

When the two factors of the second member are finite, the functions \(f\) and \(f'\) are moderated (§1, No.~2, Prop. 7), therefore the function \(f \otimes f'\) is moderated (Cor. 2 of Prop. 5); the above equality therefore reduces to the formula (10) (§1, No.~2, Prop. 7). When one of the factors of the second member has the value \(+\infty\) and the other is not zero, then the second member has the value \(+\infty\) , and the above equality follows from Prop. 6.

COROLLARY 2. --- Let f and \(f'\) be two functions with values in C or in \(\overline{R}\) , defined on T and \(T'\) , respectively, and essentially integrable (resp. integrable) for the measures \(\mu\) and \(\mu'\) , respectively. The function \(f \otimes f'\) is then essentially integrable (resp. integrable) for the measure \(\mu \otimes \mu'\) , and

\[
\iint f (t) f ^ {\prime} \left(t ^ {\prime}\right) d \mu (t) d \mu^ {\prime} \left(t ^ {\prime}\right) = \left(\int f (t) d \mu (t)\right) \left(\int f ^ {\prime} \left(t ^ {\prime}\right) d \mu^ {\prime} \left(t ^ {\prime}\right)\right).\tag{12}
\]

When \(f\) and \(f'\) are positive, \(f \otimes f'\) is measurable by Cor. 3 of Prop. 3, and the statement follows from formula (10) (resp. (11)) and the criterion for essential integrability (§1, No.~3, Prop. 9) (resp. the integrability criterion of Ch. IV, §5, No.~6, Th. 5). The general case then follows immediately.

Corollary 2 extends at once to the product of two complex measures.

